\begin{proofbox}

	\textbf{\textcolor{CProof}{思路提示}}
	从充要条件的两个方向入手：若 $f(a+x)+f(a-x)=2b$，验证图象上任意点及其对称点都在图象上；若图象关于 $(a,b)$ 对称，则恰好推出该等式。

	\medskip
	\textbf{\textcolor{CProof}{正式推导}}
	\begin{description}[style=nextline, leftmargin=2.8em, labelsep=0.5em, itemsep=0.55em]
		\item[\textbf{步骤一（必要性条件）：}] 假设 $f(x)$ 的图象关于点 $(a,b)$ 中心对称。
			\[
				\begin{aligned}
					\text{对任意 }x,\quad &(a+x,\; f(a+x)) \text{ 在图象上},\\
					&\text{它关于 }(a,b)\text{ 的对称点为 }(a-x,\; 2b - f(a+x)).
			\end{aligned}
			\]
			\par\textbf{理由：}
			由对称性，该对称点也在图象上，故 $f(a-x) = 2b - f(a+x)$，即 $f(a+x)+f(a-x)=2b$。

		\item[\textbf{步骤二（充分性条件）：}] 假设对任意 $x$ 有 $f(a+x)+f(a-x)=2b$。
			\[
				\begin{aligned}
					\text{取图象上一点 }(a+x,\; f(a+x)),\\
					\text{其关于 }(a,b)\text{ 的对称点为 }(a-x,\; 2b - f(a+x)) = (a-x,\; f(a-x)).
			\end{aligned}
			\]
			\par\textbf{理由：}
			由等式可得 $2b - f(a+x)=f(a-x)$，故该对称点也在图象上。由于定义域关于 $a$ 对称，所有点都被覆盖，因此图象关于 $(a,b)$ 中心对称。

		\item[\textbf{步骤三（等价变形）：}] 令 $g(x)=f(a+x)-b$，则
			\[
				\begin{aligned}
					g(-x) &= f(a-x)-b \\
					&= (2b - f(a+x)) - b \\
					&= - (f(a+x)-b) \\
					&= -g(x),
			\end{aligned}
			\]
			故 $g(x)$ 为奇函数。
			\par\textbf{理由：}
			这表明中心对称函数可通过平移化为奇函数，便于进一步处理。
	\end{description}

	\medskip
	\textbf{\textcolor{CProof}{结论回扣}}

	因此，$f(x)$ 的图象关于点 $(a,b)$ 中心对称 $\iff$ $f(a+x)+f(a-x)=2b$ 对定义域内所有 $x$ 成立。
\end{proofbox}
